
\documentclass{article}
\usepackage{colortbl}
\usepackage{makecell}
\usepackage{multirow}
\usepackage{supertabular}

\begin{document}

\newcounter{utterance}

\centering \large Interaction Transcript for game `cladder', experiment `full\_v1.5\_default', episode 8348 with reflexion{-}llama{-}3.1{-}8b{-}instruct.
\vspace{24pt}

{ \footnotesize  \setcounter{utterance}{1}
\setlength{\tabcolsep}{0pt}
\begin{supertabular}{c@{$\;$}|p{.15\linewidth}@{}p{.15\linewidth}p{.15\linewidth}p{.15\linewidth}p{.15\linewidth}p{.15\linewidth}}
    \# & \multicolumn{2}{c}{Player} && \multicolumn{2}{c}{Game Master} \\
    \hline

    \theutterance \stepcounter{utterance}  
    & & & \multicolumn{4}{p{0.6\linewidth}}{
        \cellcolor[rgb]{0.9,0.9,0.9}{
            \makecell[{{p{\linewidth}}}]{
                \texttt{\tiny{[P1$\langle$GM]}}
                \texttt{You are an expert in causal inference. The following question is not a typical commonsense query, but rather a meticulously designed question created by a professor specializing in causal inference, intended to assess the students' mastery of the course content.} \\
\texttt{Imagine a self{-}contained, hypothetical world with only the following conditions, and without any unmentioned factors or causal relationships: Yomx has a direct effect on cwoi and gwet. Gwet has a direct effect on xevu. Cwoi has a direct effect on xevu. The overall probability of yomx is 22\%. The probability of not yomx and xevu is 48\%. The probability of yomx and xevu is 7\%. Is the chance of xevu larger when observing yomx?} \\
\texttt{Start your answer with 'Yes' or 'No', followed by additional reasoning or evidence to support your explanation.} \\
            }
        }
    }
     \\ \\

    \theutterance \stepcounter{utterance}  
    & \multicolumn{4}{p{0.6\linewidth}}{
        \cellcolor[rgb]{0.561,0.737,0.561}{
            \makecell[{{p{\linewidth}}}]{
                \texttt{\tiny{[P1$\rangle$GM]}}
                \texttt{No.} \\
\\ 
\texttt{To determine if the chance of xevu is larger when observing yomx, we need to calculate the conditional probability of xevu given yomx (P(xevu|yomx)) and compare it to the unconditional probability of xevu (P(xevu)).} \\
\\ 
\texttt{We can use the law of total probability to calculate P(xevu|yomx):} \\
\\ 
\texttt{P(xevu|yomx) = P(xevu and yomx) / P(yomx)} \\
\\ 
\texttt{Given values:} \\
\texttt{{-} P(yomx) = 22\%} \\
\texttt{{-} P(xevu and yomx) = 7\%} \\
\\ 
\texttt{P(xevu|yomx) = 7\% / 22\% ≈ 0.3182 or 31.82\%} \\
\\ 
\texttt{Now, we need to calculate the unconditional probability of xevu (P(xevu)). We can use the law of total probability again:} \\
\\ 
\texttt{P(xevu) = P(xevu and yomx) + P(xevu and not yomx)} \\
\\ 
\texttt{Given values:} \\
\texttt{{-} P(xevu and yomx) = 7\%} \\
\texttt{{-} P(not yomx and xevu) = 48\%} \\
\\ 
\texttt{However, we are given the probability of not yomx and xevu, not the probability of xevu and not yomx. To find P(xevu), we need to calculate P(xevu and not yomx) using the fact that P(not yomx) = 1 {-} P(yomx) = 78\% and the given probability of not yomx and xevu.} \\
\\ 
\texttt{Since P(xevu and not yomx) + P(xevu and yomx) = P(xevu) and P(xevu and not yomx) + P(not yomx and xevu) = P(xevu), we can equate the two expressions:} \\
\\ 
\texttt{P(xevu and not yomx) + P(xevu and yomx) = P(xevu and not yomx) + P(not yomx and xevu)} \\
\\ 
\texttt{Subtracting P(xevu and not yomx) from both sides:} \\
\\ 
\texttt{P(xevu and yomx) = P(not yomx and xevu)} \\
\\ 
\texttt{Given values:} \\
\texttt{{-} P(xevu and yomx) = 7\%} \\
\texttt{{-} P(not yomx and xevu) = 48\%} \\
\\ 
\texttt{However, the given value of 48\% is the probability of not yomx and xevu, not the probability of xevu and not yomx. We can't directly use this value to calculate P(xevu).} \\
\\ 
\texttt{To find P(xevu), we need to use the fact that P(xevu) = P(xevu and yomx) + P(xevu and not yomx)} \\
            }
        }
    }
    & & \\ \\

    \theutterance \stepcounter{utterance}  
    & & & \multicolumn{2}{p{0.3\linewidth}}{
        \cellcolor[rgb]{0.9,0.9,0.9}{
            \makecell[{{p{\linewidth}}}]{
                \texttt{\tiny{[GM$|$GM]}}
                \texttt{No} \\
            }
        }
    }
    & & \\ \\

    \theutterance \stepcounter{utterance}  
    & & & \multicolumn{2}{p{0.3\linewidth}}{
        \cellcolor[rgb]{0.9,0.9,0.9}{
            \makecell[{{p{\linewidth}}}]{
                \texttt{\tiny{[GM$|$GM]}}
                \texttt{no} \\
            }
        }
    }
    & & \\ \\

    \theutterance \stepcounter{utterance}  
    & & & \multicolumn{2}{p{0.3\linewidth}}{
        \cellcolor[rgb]{0.9,0.9,0.9}{
            \makecell[{{p{\linewidth}}}]{
                \texttt{\tiny{[GM$|$GM]}}
                \texttt{game\_result = WIN} \\
            }
        }
    }
    & & \\ \\

\end{supertabular}
}

\end{document}
